\documentclass[11pt,letterpaper]{article}
\usepackage[utf8]{inputenc}
\usepackage[spanish,es-tabla]{babel}
\usepackage{amsmath}
\usepackage{amsfonts}
\usepackage{graphicx}
\usepackage{booktabs}
\usepackage{geometry}
\usepackage{hyperref}
\usepackage{cite}

\geometry{top=2.5cm, bottom=2.5cm, left=2.5cm, right=2.5cm}

\title{
    \small \textbf{UNIVERSIDAD DE CONCEPCIÓN}\\
    Facultad de Ingeniería\\
    Departamento de Ingeniería Informática y Ciencias de la Computación\\[1cm]
    \large Proyecto Semestral - Estructuras de Datos\\[0.5cm]
    \huge \textbf{Análisis Comparativo de Estructuras de Datos para el Procesamiento de Volúmenes de Datos}}

\author{
    \textbf{Integrantes:}\\
    Nicolás Ricciardi \\
    Enzo Levancini \\
    Máximo Beltrán
}
\date{Julio 2026}

\begin{document}

\maketitle
\thispagestyle{empty}
\newpage

\tableofcontents
\newpage

% ==========================================
% 1. INTRODUCCIÓN (Máximo 1 página)
% ==========================================
\section{Introducción}
\label{sec:introduccion}
% Breve reseña de todo el trabajo realizado. 
% Descripción a alto nivel del problema planteado y estructuras de datos utilizadas (implementadas y de bibliotecas).
% Herramientas, fuentes de datos y conclusiones preliminares.

Aquí va la reseña general del trabajo. Se debe describir a alto nivel el problema del procesamiento y conteo de tweets por usuario, justificando la necesidad de evaluar diferentes estructuras de datos tanto de bibliotecas estándar como implementadas en el curso. 

Mencionar brevemente las herramientas utilizadas (ej. C++17, MinGW-w64), la fuente de datos (dataset de tweets) y adelantar las conclusiones preliminares clave respecto a qué estructura y qué tipo de clave resultó ser más eficiente.


% ==========================================
% 2. DESCRIPCIÓN DE ESTRUCTURAS DE DATOS (Máximo 1 página)
% ==========================================
\section{Descripción de las Estructuras de Datos Comparadas}
\label{sec:estructuras}
% Repositorio de código
Los códigos fuentes documentados correspondientes a las implementaciones analizadas se encuentran disponibles de manera adjunta y en el siguiente repositorio de acceso público: \url{https://github.com/tu-usuario/tu-repositorio}. El desarrollo sigue el estándar de codificación adoptado para el curso.

\subsection{Estructura de Datos 1 (ej. unordered\_map / Hash Table)}
Descripción general de la primera estructura. Indicar detalladamente los costos teóricos de inserción y búsqueda en notación Big-O ($O(1)$ promedio, $O(n)$ peor caso). Especificar los parámetros de configuración utilizados (función hash implementada, número inicial de buckets, factor de carga máximo). \textit{Restricción: Máximo un párrafo por estructura, incluyendo referencias cruzadas o links si aplica \cite{cormen2009}.}% Mantener conciso.

\subsection{Estructura de Datos 2 (ej. map / AVL / Red-Black Tree)}
Descripción general de la segunda estructura. Indicar los costos de inserción y búsqueda ($O(\log n)$). Especificar los parámetros utilizados en la experimentación. \textit{Restricción: Máximo un párrafo por estructura.}

\subsection{Estructura de Datos 3 (ej. Estructura Propia / Custom Hash)}
Descripción general de la tercera estructura. Indicar costos de inserción y búsqueda, detallando el manejo de colisiones (ej. encadenamiento o direccionamiento abierto) y los parámetros específicos elegidos. \textit{Restricción: Máximo un párrafo por estructura.}


% ==========================================
% 3. DATASET Y AMBIENTE EXPERIMENTAL (Máximo 1 página)
% ==========================================
\section{Descripción del Dataset y del Ambiente de Experimentación}
\label{sec:ambiente}

\subsection{Lectura y Modificación del Dataset}
Describir detalladamente los métodos de lectura implementados para procesar el dataset de tweets de forma eficiente (ej. parsing de líneas, manejo de buffers de entrada). Si fue necesario realizar modificaciones o limpieza previa del dataset (filtrado de nulos, corrección de codificación), se deben especificar y justificar aquí estructuralmente.

\subsection{Ambiente de Hardware y Software}
Los experimentos se llevaron a cabo en una máquina con las siguientes especificaciones técnicas detalladas en la Tabla \ref{tab:hardware}:

\begin{table}[h!]
\centering
\caption{Especificaciones del Entorno de Pruebas}
\label{tab:hardware}
\begin{tabular}{ll}
\toprule
\textbf{Componente} & \textbf{Detalle} \\
\midrule
Procesador & [Modelo exacto, ej. AMD Ryzen 5 5600H] \\
Memoria RAM & [Capacidad y velocidad, ej. 16 GB DDR4 3200 MHz] \\
Memoria Caché & L1: [e.g., 384 KB], L2: [e.g., 3 MB], L3: [e.g., 16 MB] \\
Sistema Operativo & [ej. Windows 11 Home / Ubuntu 22.04 LTS] \\
Compilador & [ej. MinGW-w64 GCC 14.x] con banderas de optimización \texttt{-O2} \\
\bottomrule
\end{tabular}
\end{table}


% ==========================================
% 4. RESULTADOS EXPERIMENTALES (Máximo 4 páginas)
% ==========================================
\section{Resultados Experimentales}
\label{sec:resultados}

\subsection{Tamaño de las Estructuras de Datos en Memoria}
Se evaluó el consumo espacial de cada estructura cargada con el total de los datos. Los resultados expresados en KB o MB se resumen en la Tabla \ref{tab:memoria}.

\begin{table}[h!]
\centering
\caption{Comparativa de Consumo de Memoria}
\label{tab:memoria}
\begin{tabular}{lcc}
\toprule
\textbf{Estructura de Datos} & \textbf{Tamaño con \texttt{user\_id} (MB)} & \textbf{Tamaño con \texttt{user\_screen\_name} (MB)} \\
\midrule
Estructura 1 & & \\
Estructura 2 & & \\
Estructura 3 & & \\
\bottomrule
\end{tabular}
\end{table}

\subsection{Rendimiento en la Creación y Escalabilidad}
Para medir la escalabilidad en la creación de las estructuras (conteo de tweets por usuario), se registró de forma incremental el tiempo de ejecución en milisegundos a medida que aumentaba el volumen de tweets procesados (desde 10.000 hasta 180.000). 

% Ejemplo de cómo incluir gráficos (deben estar en la misma carpeta)
% \begin{figure}[h!]
% \centering
% \includegraphics[width=0.7\textwidth]{grafico_escalabilidad.png}
% \caption{Tiempo de ejecución en la creación versus cantidad de tweets procesados.}
% \label{fig:escalabilidad}
% \end{figure}

\subsection{Análisis del Impacto del Tipo de Clave}
Se contrastó el rendimiento de las búsquedas e inserciones utilizando \texttt{user\_id} (numérico/entero) frente a \texttt{user\_screen\_name} (cadena de texto / \texttt{std::string}). Se entrega evidencia experimental explícita de cuál clave resulta más eficiente para el cálculo de las funciones de dispersión o comparaciones de nodos en cada tipo de estructura.


% ==========================================
% 5. CONCLUSIONES (Máximo 1 página)
% ==========================================
\section{Conclusiones}
\label{sec:conclusiones}
% Describir los principales descubrimientos obtenidos. 
% Hacer hincapié en las estructuras de datos más rápidas para cada tipo de experimento y en la comparación de los dos tipos de claves evaluadas.

Se detallan los hallazgos principales derivados del estudio empírico. Se destaca de forma clara y justificada qué estructura de datos demostró el mejor tiempo de respuesta ante la escalabilidad de carga masiva de tweets.

Asimismo, se discute el impacto del tipo de clave evaluada, concluyendo sobre la ventaja de desempeño que supone el uso de \texttt{user\_id} sobre \texttt{user\_screen\_name} debido a costos reducidos de comparación lógica y cálculo de hashes, cerrando con la relación general costo-beneficio de tiempo versus memoria.

\newpage
% ==========================================
% REFERENCIAS
% ==========================================
\begin{thebibliography}{9}
\bibitem{cormen2009}
T. H. Cormen, C. E. Leiserson, R. L. Rivest, y C. Stein, \textit{Introduction to Algorithms}, 3rd ed. MIT Press, 2009.
\end{thebibliography}

\end{document}